La constante de Catalan apparaît habituellement comme une série alternée :\par

\[
G
=
1-\frac1{3^2}+\frac1{5^2}-\frac1{7^2}+\cdots.
\]

Cette formule donne sa valeur et la généalogie HMT explique pourquoi cette alternance particulière existe.\par

Dans HMT, l’origine se trouve dans le quart de torsion.\par

Il existe un opérateur interne \(J\) tel que\par

\[
J^2=-I.
\]

La racine de \(-1\) émerge de deux applications successives d’un quart de tour. La structure complexe est la mémoire algébrique d’une rotation interne.\par

Cette rotation engendre un caractère d’ordre quatre. Lorsqu’il est publié sur les entiers positifs, il produit exactement le motif résiduel qui alterne entre \(+1\), \(0\), \(-1\), \(0\).\par

La constante de Catalan apparaît alors comme une trace spectrale :\par

\[
G
=
L(2,\chi_{-4})
=
\operatorname{Tr}(N^{-2}Q_4).
\]

Cela change sa signification. Catalan est simultanément la somme d’une série particulière et la réponse spectrale des entiers à un quart de torsion.\par

Chaque terme de la série enregistre la manière dont la rotation interne agit sur une classe résiduelle. L’alternance est l’ombre arithmétique de la rotation.\par

La structure triadique possède son propre caractère modulo \(3\). En le combinant au caractère d’ordre quatre, le théorème chinois des restes produit un caractère modulo \(12\) :\par

\[
\chi_{12}
=
\chi_{-3}\chi_{-4}.
\]

Le module \(12\) apparaît comme l’espace commun minimal dans lequel l’orientation triadique et le quart de torsion peuvent coexister.\par

La dodécaphase est, en ce sens, une réconciliation de \(3\) et de \(4\).\par

Du caractère combiné émergent\par

\[
L(1,\chi_{12})
=
\frac{\log(2+\sqrt3)}{\sqrt3},
\]

y\par

\[
L(2,\chi_{12})
=
\frac{\pi^2}{6\sqrt3}.
\]

La cantidad\par

\[
2+\sqrt3
\]

est l’unité fondamentale du corps réel \(\mathbb Q(\sqrt3)\). Son logarithme représente un déplacement hyperbolique. De nouveau, une structure finie de résidus ouvre une dynamique multiplicative infinie.\par

Ici se rencontrent rotation, torsion, caractère arithmétique, unité de Pell et déplacement hyperbolique.\par

La constante d’Euler–Kronecker du corps \(\mathbb Q(\sqrt3)\) prolonge cette structure. Elle mesure le terme régulier qui demeure après la séparation du pôle principal de la fonction zêta du corps. Dans la généalogie HMT, ce terme constitue l’invariant régulier de mémoire arithmétique fine qui subsiste après le retrait de la divergence universelle.\par

La famille Bendersky–Glaisher prolonge le mécanisme en une tour. Les dérivées spectrales de zêta engendrent cette tour. Apéry et d’autres constantes associées à des valeurs impaires s’organisent comme membres d’une même famille fonctionnelle.\par

L’apport nouveau réside dans la généalogie qui réunit les constantes classiques :\par

\begin{itemize}[leftmargin=*,itemsep=0.35em]
\item le quart de torsion produit la structure complexe ;
\item le résidu triadique produit la structure modulo \(3\) ;
\item leur composition produit le caractère modulo \(12\) ;
\item le caractère engendre des valeurs \(L\) ;
\item les valeurs \(L\) ouvrent le corps quadratique et ses constantes spectrales.
\end{itemize}

La torsion géométrique est devenue arithmétique en conservant intégralement son orientation.\par
